\documentclass[12pt]{article}
\usepackage[margin=0.6in,top=1.15in,bottom=0.55in]{geometry}
\usepackage{lmodern}
\usepackage{microtype}
\usepackage{titlesec}
\usepackage{enumitem}
\usepackage[hidelinks]{hyperref}
\usepackage{../../latex/grader-worksheet}

\setlength{\parindent}{0pt}
\setlength{\parskip}{0.25em}
\titleformat{\section}{\large\bfseries\scshape}{}{0pt}{}

\GraderSetup{assignment-id=U1L3LE,template-revision=1}

\begin{document}

\begin{center}
{\LARGE\bfseries Week 3 Lit Example Worksheet}\par\vspace{0.05em}
{\large\scshape Macbeth: Prophecies as Propositions}\par\vspace{0.12em}
\rule{0.78\textwidth}{0.5pt}
\end{center}

\textbf{Purpose:} Apply Week 3's Whately method to Macbeth. Decide whether a line is a proposition, translate clear claims into standard categorical form, and identify quality, quantity, and distribution.

\textbf{Directions:} Use the Lit Example Reader. Quote only briefly. Your task is logical translation, not proving whether the witches are trustworthy.

\section*{Part 1: Proposition or Not?}

\textbf{Question 1.1}\\[-0.15em]
The First Apparition says, ``Beware Macduff.'' Is this a proposition in the strict Week 3 sense? Explain why or why not. Then write one proposition Macbeth might infer from the warning.

\GraderAnswerBox{Part 1 Q1}{2.15in}

\textbf{Question 1.2}\\[-0.15em]
The Second Apparition says, ``None of woman born / Shall harm Macbeth.'' Is this a proposition? Identify its subject and predicate in ordinary language.

\GraderAnswerBox{Part 1 Q2}{2.15in}

\newpage

\section*{Part 2: Standard Form}

\textbf{Question 2.1}\\[-0.15em]
Translate ``None of woman born shall harm Macbeth'' into a clear E-form proposition: ``No S are P.'' Then state its quality and quantity.

\GraderAnswerBox{Part 2 Q1}{2.15in}

\textbf{Question 2.2}\\[-0.15em]
Macbeth treats the apparition's claim as if it proves, ``All men I will meet in battle are people who cannot harm me.'' Name the standard form, quality, and quantity of Macbeth's translated claim. Then explain why this may go beyond the exact wording.

\GraderAnswerBox{Part 2 Q2}{2.45in}

\newpage

\section*{Part 3: Translation and Distribution}

\textbf{Question 3.1}\\[-0.15em]
Translate this claim into standard form: ``Only someone not born in the ordinary way can harm Macbeth.'' State the subject and predicate you chose.

\GraderAnswerBox{Part 3 Q1}{2.35in}

\textbf{Question 3.2}\\[-0.15em]
For your standard-form proposition in Part 3 Q1, identify which term is distributed and which term is undistributed. Briefly explain.

\GraderAnswerBox{Part 3 Q2}{2.35in}

\newpage

\section*{Part 4: Testing Macbeth's Confidence}

\textbf{Question 4.1}\\[-0.15em]
Macduff says he was ``from his mother's womb / Untimely ripp'd.'' Explain how this fact affects the subject class in the proposition ``No persons born of woman are persons who will harm Macbeth.'' Does Macduff clearly fall inside or outside that class? Explain the logical issue.

\GraderAnswerBox{Part 4 Q1}{2.55in}

\textbf{Question 4.2}\\[-0.15em]
Put Whately's Week 3 question to Macbeth: before relying on a prophecy, what exact proposition must he translate and test? Write the proposition, identify its form if possible, and state one danger of treating unclear dramatic language as a settled logical guarantee.

\GraderAnswerBox{Part 4 Q2}{2.55in}

\end{document}
